%\documentclass[12pt,letterpaper]{article}



%%%
%% Required packages:
%%
\usepackage{etex}
\usepackage{amsmath}
\usepackage{amssymb}
\usepackage{amstext}
\usepackage{amsthm}
\usepackage{fancyhdr,fancybox}
\usepackage{mathrsfs} % need for \mathscr command
\usepackage{multicol}
\usepackage[normalem]{ulem}
%\usepackage{pictex,pictexplus}
% \usepackage{pictexwd}
\usepackage{rotating}
\usepackage{url}
\usepackage{xfrac}
\usepackage{booktabs}
\usepackage{color,soul}
\usepackage{comment}

\usepackage{hyperref}
\hypersetup{
    colorlinks=true,     % false: boxed links; true: colored links
    linkcolor=blue,      % color of internal links
    citecolor=blue,      % color of links to bibliography
    filecolor=blue,      % color of file links
    urlcolor=blue        % color of external links
}


\usepackage{tikz}
\usetikzlibrary{backgrounds,calc}

%%
%% Common formatting commands
%%
\setlength{\parindent}{0in}
\setlength{\textwidth}{7in}
\setlength{\evensidemargin}{-0.25in}
\setlength{\oddsidemargin}{-0.25in}
\setlength{\parskip}{.5\baselineskip}
\setlength{\topmargin}{-0.5in}
\setlength{\textheight}{9in}

%               Problem and Part
\newcounter{problemnumber}
\newcounter{partnumber}
\newcounter{subpartnumber}
\newcommand{\Problem}{%
  \refstepcounter{problemnumber}%
  \setcounter{partnumber}{0}%
  \setcounter{subpartnumber}{0}%
  \item[\makebox{\hfil\textbf{\theproblemnumber.}\hfil}]%
  }
\newcommand{\InProblem}[2][1.6]{%
  \refstepcounter{problemnumber}%
  \setcounter{partnumber}{0}%
  \setcounter{subpartnumber}{0}%
  \textbf{\theproblemnumber.}\ \ \parbox[t]{#1in}{#2}%
}
\newcommand{\InSmallProblem}[1]{\InProblem[1.05]{#1}}
\newcommand{\NoProblem}[1]{%
  \mbox{\hfil\textbf{#1}\hfil}%
  }
\newcommand{\UnProblem}{%
  \refstepcounter{problemnumber}%
  \setcounter{partnumber}{0}%
  \setcounter{subpartnumber}{0}%
  \mbox{\hfil\textbf{\theproblemnumber.}\hfil}%
  }
\newcommand{\InFlexProblem}[1]{%
  \refstepcounter{problemnumber}%
  \setcounter{partnumber}{0}%
  \setcounter{subpartnumber}{0}%
  \textbf{\theproblemnumber.}\ \ #1%
}
%%  Part:
\newcommand{\Part}{%
  \renewcommand{\thepartnumber}{(\alph{partnumber})}%
  \refstepcounter{partnumber}
  \setcounter{subpartnumber}{0}%
  \item[(\alph{partnumber})]%
}
\newcommand{\InPart}[2][1.75]{%
  \renewcommand{\thepartnumber}{(\alph{partnumber})}%
  \refstepcounter{partnumber}%
  \setcounter{subpartnumber}{0}%
  (\alph{partnumber})\ \ \parbox[t]{#1in}{#2}%
}
\newcommand{\UnPart}{%
  \refstepcounter{partnumber}%
  \renewcommand{\thepartnumber}{(\alph{partnumber})}%
  \setcounter{subpartnumber}{0}%
  (\alph{partnumber})%
}
\newcommand{\InLargePart}[1]{\InPart[3.05]{#1}}
\newcommand{\InFullPart}[1]{\InPart[6.2]{#1}}
\newcommand{\InSmallPart}[1]{\InPart[1.05]{#1}}
%% SubPart:
\newcommand{\SubPart}{%
  \refstepcounter{subpartnumber}%
  \item[(\roman{subpartnumber})]%
  }
\newcommand{\InSubPart}[2][2.25]{%
  \stepcounter{subpartnumber}%
  (\roman{subpartnumber})\ \ \parbox[t]{#1in}{#2}%
}
\newcommand{\InSmallSubPart}[1]{\InSubPart[1.5]{#1}}
\newcommand{\InTinySubPart}[1]{%
  \stepcounter{subpartnumber}%
  (\roman{subpartnumber})\ \ \makebox{#1}%
}
\newcommand{\InMySubPart}[2][2.25]{%
  \stepcounter{subpartnumber}%
  \parbox[t]{#1in}{\makebox[0.3in][l]{(\roman{subpartnumber})}\ \ {#2}}%
}
\newcommand{\TrueFalse}{\textbf{T}\ \ \textbf{F}\ \ }
\newcommand{\TFPart}[1]{% 
  \stepcounter{partnumber}%
  \setcounter{subpartnumber}{0}%
  \item[(\alph{partnumber})] \TrueFalse\parbox[t]{5.5in}{#1}}

\newcommand{\Points}[1]{(#1 points) \ }
\newcommand{\Point}[1]{(#1 point) \ }


%% For Answers:
\newcounter{answernumber}
\newcommand{\Answer}{%
  \refstepcounter{answernumber}%
  \setcounter{partnumber}{0}%
  \setcounter{subpartnumber}{0}%
  \item[\makebox{\hfil\textbf{\theanswernumber.}\hfil}]%
  }


% Setting up the headers for the first page
%\chead{} % will default to what is typed in the file.
\fancyhead[L]{\textsc{Math 22a}}
\fancyhead[R]{\textsc{Fall 2024}}
\fancyfoot[R]{
	\vspace*{\fill}\footnotesize\textcopyright{} \emph{Harvard Math 22a}	
}
\renewcommand{\headrulewidth}{0pt}

%\fancyhead[C]{}
%	\chead{}	
%	\lhead{}
%	\rhead{}
%	\cfoot{}


% Putting copyright on the footer of each page
%\fancypagestyle{secondpage}{
%	\fancyhead[C]{TESTing again} % change this to be blank
%		%\chead{}	
%	\fancyhead[L]{bears} % change this to be blank
%	\fancyhead[R]{dogs} % change this to be blank
%	\fancyfoot[R]{ %keeps the footer the same
%		\vspace*{\fill}\footnotesize\textcopyright{} \emph{Harvard Math 22a}	
%	}
%}
%\renewcommand{\headrulewidth}{0pt}
%\pagestyle{fancy}
\pagestyle{empty}


%\lhead{}
%\rhead{}
%\rfoot{\vspace*{\fill}\footnotesize\textcopyright{} \emph{Harvard Math 22a}}
%\renewcommand{\headrulewidth}{0pt}
%\pagestyle{fancy}




%%
%% Common shortcut commands:
%%
\newcommand{\ds}{\displaystyle}
\newcommand{\sgn}{\operatorname{sgn}}
\renewcommand{\Pr}{\operatorname{P}}
\newcommand{\CPr}[3][{}]{\Pr_{\text{#1}}\left(#2\,|\,#3\right)}

\newcommand{\C}{\mathbb{C}}
\newcommand{\R}{\mathbb{R}}
\newcommand{\Z}{\mathbb{Z}}
\newcommand{\N}{\mathbb{N}}
\newcommand{\Q}{\mathbb{Q}}
\newcommand{\D}{\mathbb{D}}

\newcommand{\rref}{\operatorname{rref}}
\newcommand{\im}{\operatorname{im}}
\newcommand{\rank}{\operatorname{rank}}
\newcommand{\diag}{\operatorname{diag}}
\newcommand{\Span}{\operatorname{span}}
%\renewcommand{\vec}[1]{\mathbf{#1}}
\newcommand{\charpoly}{\mathscr{P}}
\newcommand{\nicefrac}[2]{\sfrac{#1}{#2}} % using xfrac package
\newcommand{\topology}{\mathcal{T}}
\newcommand{\Inn}{\operatorname{Inn}}

\newcommand{\blankbox}[2]{\fbox{\rule{#1}{0in}\rule{0in}{#2}}}

\newenvironment{myindentpar}[1]%
 {\begin{list}{}%
         {\setlength{\leftmargin}{#1}
          \setlength{\rightmargin}{\leftmargin}}%
         \item[]%
 }
 {\end{list}}
\newtheorem{theorem}{Theorem}
\newtheorem*{NStheorem}{Theorem}
\newtheorem{cor}[theorem]{Corollary}
\newtheorem*{claim}{Claim}
\newtheorem{defn}{Definition}

\newenvironment{amatrix}[1]{%
  \left(\begin{array}{@{}*{#1}{c}|c@{}}
}{%
  \end{array}\right)
}

%--------grstep
% For denoting a Gauss' reduction step.
% Use as: \grstep{\rho_1+\rho_3} or \grstep[2\rho_5 \\ 3\rho_6]{\rho_1+\rho_3}
\newcommand{\grstep}[2][\relax]{%
   \ensuremath{\mathrel{
       {\mathop{\longrightarrow}\limits^{#2\mathstrut}_{
                                     \begin{subarray}{l} #1 \end{subarray}}}}}}
\newcommand{\swap}{\leftrightarrow}




\section{{Linear Independence,  $\vec\S$1.7}}% 
\chead{\Large\textbf{Linear Independence,  $\vec\S$1.7}}% 

%\begin{document}
\thispagestyle{fancy}

Recall from last class:

\begin{itemize}

\item A linear system is called {\bf homogeneous} if it can be written as $A\vec{x}=\vec{0}$ where $A$ is an $m \times n$ matrix, $\vec{x} \in \mathbb{R}^n$, and $\vec{0}$ is the vector of $m$ zeros. 
\item {\bf Theorem:} If $A\vec{x} =\vec{b}$ is consistent for a given $\vec{b}$, and $\vec{p}$ is a solution, then the solution set of $A\vec{x}=\vec{b}$ is the set of vectors of the form $\vec{p}+\vec{v}_h$ where $\vec{v}_h$ is a solution of $A\vec{x} =\vec{0}$.
\end{itemize}

\newpage

\begin{defn}
A set of vectors $\{ \vec{v}_1, \dots, \vec{v}_p\}$ in $\mathbb{R}^n$ is said to be {\bf linearly independent} if $$x_1 \vec{v}_1 +x_2 \vec{v}_2 +\dots +x_p \vec{v}_p = \vec{0}$$ has only the trivial solution. The set is called {\bf linearly dependent} if the equation has non-trivial solutions and the equation is called the {\bf dependence relation}.
\end{defn}

\begin{enumerate}

\Problem Let $\vec{v}_1=\begin{bmatrix} 1\\2\\3\\ \end{bmatrix}$, $\vec{v}_2=\begin{bmatrix} 4\\5\\6\\ \end{bmatrix}$, and $\vec{v}_3=\begin{bmatrix} 2\\1\\0\\ \end{bmatrix}$. Is $\{ \vec{v}_1, \vec{v}_2,\vec{v}_3 \}$ linearly independent?

\vfill

\vfill

\vfill

Let $A= [\vec{a}_1 \; \dots \; \vec{a}_n ]$, then $\{ \vec{a}_1 \; \dots \; \vec{a}_n \}$ are linearly independent if and only if $A \vec{x} = \vec{0}$ has only the trivial solution.

\Problem Is a set with only one vector always linearly independent?

\vfill

\Problem If $\vec{0} \in \{ \vec{v}_1, \dots, \vec{v}_p \}$, then is the set linearly dependent?

\vfill

\Problem If $\vec{0} \in \text{Span}\{ \vec{v}_1, \dots, \vec{v}_p \}$, then the vectors are linearly independent?

\vfill

\newpage

\Problem Let $\vec{v}_1=\begin{bmatrix} 1\\2\\ \end{bmatrix}$, $\vec{v}_2=\begin{bmatrix} 3\\4\\ \end{bmatrix}$, and $\vec{v}_3=\begin{bmatrix} 5\\6\\ \end{bmatrix}$. Are the vectors linearly independent?

\vfill

{\bf Theorem 7:} A set $S=\{ \vec{v}_1, \dots, \vec{v}_p \}$ of two or more vectors is linearly dependent if and only if one of the vectors of $S$ is a linearly combination of the others.

Moreover if $\vec{v}_1\neq \vec{0}$, then some $\vec{v}_j$ for $j>1$ is a linear combination of $\vec{v}_1, \dots, \vec{v}_{j-1}$.

\Problem Prove Theorem 7.

\vfill

\newpage

\Problem Let $\vec{u}=\begin{bmatrix} 1\\ 0\\ 1\\ \end{bmatrix}$ and $\vec{v} = \begin{bmatrix} 2 \\ 0 \\ 8\\ \end{bmatrix}$
\Part Describe $\text{Span}\{\vec{u}, \vec{v}\}$

\vfill

\Part Prove $\vec{w} \in \text{Span}\{\vec{u}, \vec{v}\}$ if and only if $\{ \vec{u}, \vec{v}, \vec{w} \}$ is linearly dependent.

\vfill
\vfill

{\bf Theorem 8:} Any set $\{ \vec{v}_1, \dots, \vec{v}_p \}$ in $\mathbb{R}^n$ is linearly dependent if $p>n$.

\Problem Prove the theorem.

\vfill
\vfill
\vfill 
\end{enumerate}

% sols start below
%\end{document}
\begin{comment}

\newpage
\chead{\Large\textbf{Linear Independence -- Answers / Solutions}}%
\lhead{}
\rhead{}
\thispagestyle{fancy}

\begin{enumerate}
  \Answer In order to determine if the vectors are linearly independent, we need to solve the vector equation $c_1 \vec{v}_1+c_2 \vec{v}_2+c_3 \vec{v}_3=\vec{0}$. We instead solve the corresponding augmented matrix. Row-reducing the augmented matrix, we get
  $$\begin{amatrix}{3}
  1 & 4 & 2 & 0\\
  2 & 5 & 1 & 0\\
  3 & 6 & 0 & 0\\
  \end{amatrix}\sim
  \begin{amatrix}{3}
  1 & 0 & -2 & 0\\
0 & 1 & 1 & 0\\
0 & 0 & 0 & 0\\  
\end{amatrix}.$$
Thus we get $c_1-2c_3=0$ and $c_2+c_3=0$, and hence the solution set is given by $$\left\{ \vec{c} \in \mathbb{R}^3 : c_1=2c_3, c_2=-c_3\right\}=
\left\{ \begin{bmatrix}
2c_3\\
-c_3\\
c_3\\
\end{bmatrix} : c_3 \in \mathbb{R} \right\}=
\left\{ c_3 \begin{bmatrix}
2\\
-1\\
1\\
\end{bmatrix} : c_3 \in \mathbb{R} \right\}=\text{Span}\left\{\begin{bmatrix}
2\\
-1\\
1\\
\end{bmatrix} \right\}.$$
Thus the vector equation has non-trivial solutions, and the vectors are linearly dependent. From the solution set we can see how the vectors are related to one another using a dependence relation. In this case, a dependence relation is $$2\vec{v}_1-\vec{v}_2+\vec{v}_3=\vec{0}.$$

\Answer No. Since the vector equation $c \cdot \vec{0}= \vec{0}$ has non-trivial solutions, we get $\vec{0}$ is linearly dependent as a single vector. On the other hand, if $\vec{v}$ is any non-zero vector, then it is linearly independent. 

\Answer Yes. We prove the claim. Since $\vec{0} \in \{ \vec{v}_1, \dots, \vec{v}_p \}$, we know $\vec{v}_k=\vec{0}$ for some $k$. Thus $$0 \cdot \vec{v}_1+ \dots + 0 \cdot \vec{v}_{k-1}+1\cdot \vec{v}_k+ 0\cdot \vec{v}_{k+1}+\dots +0\cdot \vec{v}_p =\vec{0},$$ which gives a non-trivial solution to $x_1\vec{v}_1+\dots +x_n \vec{v}_p=\vec{0}$. Thus the vectors are linearly dependent.

\Answer No. $\vec{v} \in \text{Span}\{ \vec{v}_1, \dots, \vec{v}_p \}$ whether or not the vectors are linearly independent.

\Answer Form the augmented matrix and row-reduce to get 
  $$\begin{amatrix}{3}
1 & 3 & 5 & 0\\
2 & 4 & 6 & 0\\
\end{amatrix}\sim
\begin{amatrix}{3}
1 & 0 & -1 & 0\\
0 & 1 & 2 & 0\\  
\end{amatrix}.$$
Thus the vector equation has non-trivial solutions, and the vectors are linearly dependent. 
\Answer
\begin{proof}
$(\implies)$. Suppose $S$ is linearly dependent. By definition of linear dependence, there exists $c_1,\dots, c_p$, not all zero, such that $$c_1\vec{v}_1+\dots+ c_p \vec{v}_p =\vec{0}.$$ Since not all $c_k$ are zero, we know there exists some $c_k\neq0$. We solve the vector equation for $\vec{v}_k$ to get $$\vec{v}_k=-\frac{c_1}{c_k}\vec{v}_1-\dots--\frac{c_{k-1}}{c_k}\vec{v}_{k-1}-\frac{c_{k+1}}{c_k}\vec{v}_{k+1}-\dots -\frac{c_p}{c_k}\vec{v}_p.$$ Thus $\vec{v}_k$ can be written as a linear combination of the other vectors.

$(\impliedby)$. Suppose one of the vectors can be written as a linear combination of the others. We can assume $\vec{v}_p$ can be written in terms of the others (if not we can reorder the vectors). Thus there exists $c_1,\dots, c_{p-1}$ such that $$\vec{v}_p=c_1 \vec{v}_1 +\dots+ c_{p-1} \vec{v}_{p-1}.$$ Rearranging we get $$\vec{0}=c_1 \vec{v}_1 +\dots+ c_{p-1} \vec{v}_{p-1}-\vec{v}_p.$$ Thus we have a non-trivial solution to $x_1 \vec{v}_1 +\dots+ x_{p} \vec{v}_{p}=\vec{0}$, and thus the vectors are linearly dependent.
\end{proof}

\Answer 
\Part In this problem we have $\text{Span}\{\vec{u}, \vec{v}\}$ is the $xz$-plane in $\mathbb{R}^3$.
\Part 
\begin{proof}
$(\implies)$. Suppose $\vec{w} \in \text{Span}\{ \vec{u}, \vec{v}\}$, then $\vec{w}$ can be written as a linear combination of $\vec{u}$ and $\vec{v}$. Thus, by Theorem 7, we get $\{\vec{u},\vec{v},\vec{w}\}$ is linearly dependent.

$(\impliedby)$. Since $\vec{u}\neq 0$ and $\vec{v}$ is not parallel with $\vec{u}$, we can use Theorem 7 to conclude that $\vec{w}$ is a linear combination of $\vec{u}$ and $\vec{v}$, and hence $\vec{w}$ is in the span of $\vec{u}$ and $\vec{v}$.
\end{proof}
\Answer
\begin{proof}
Let $A=[\vec{v}_1 \; \dots \; \vec{v}_p]$ with $\vec{v}\in \mathbb{R}^n$ and $p>n$. Thus $A\vec{x}=\vec{0}$ gives a linear system of equations with $n$ equations and $p$ variables. Thus we can have at most $n$ pivots. Since $p>n$, we have at least one non-pivot column, and hence we have a free variable. Thus $A\vec{x}=\vec{0}$ has non-trivial solutions, and thus the columns are linearly dependent. 
\end{proof}
\end{enumerate}

\end{comment}

%\end{document}


%%% Local ables: 
%%% mode: latex
%%% TeX-master: t
%%% End: 
